\section{Himpunan FIRST dan FOLLOW}
\label{sec:ch10-first-follow}

\subsection{Definisi FIRST}

\begin{definisi}[FIRST]
\(\text{FIRST}(\alpha)\) adalah himpunan terminal yang dapat muncul
sebagai simbol pertama pada string yang diturunkan dari \(\alpha\).
Jika \(\alpha\) dapat menurunkan \(\varepsilon\), maka
\(\varepsilon \in \text{FIRST}(\alpha)\).
\end{definisi}

Aturan perhitungan:
\begin{enumerate}[label=\arabic*.,leftmargin=*]
  \item jika \(X\) terminal, \(\text{FIRST}(X) = \{X\}\);
  \item jika \(X \to \varepsilon\), tambahkan \(\varepsilon\);
  \item jika \(X \to Y_1Y_2\cdots Y_k\), tambahkan
        \(\text{FIRST}(Y_1) \setminus \{\varepsilon\}\); jika
        \(Y_1\) nullable, lanjut ke \(Y_2\), dan seterusnya;
  \item ulangi sampai tidak ada perubahan.
\end{enumerate}

\subsection{Definisi FOLLOW}

\begin{definisi}[FOLLOW]
\(\text{FOLLOW}(A)\) untuk nonterminal \(A\) adalah himpunan terminal
yang dapat muncul tepat setelah \(A\) pada suatu bentuk kalimat.
Simbol akhir masukan \(\$\) ditambahkan ke FOLLOW simbol awal.
\end{definisi}

Aturan perhitungan:
\begin{enumerate}[label=\arabic*.,leftmargin=*]
  \item tambahkan \(\$\) ke \(\text{FOLLOW}(S)\);
  \item untuk aturan \(A \to \alpha B \beta\), tambahkan
        \(\text{FIRST}(\beta) \setminus \{\varepsilon\}\) ke
        \(\text{FOLLOW}(B)\);
  \item jika \(\beta\) nullable atau kosong, tambahkan
        \(\text{FOLLOW}(A)\) ke \(\text{FOLLOW}(B)\);
  \item ulangi sampai stabil.
\end{enumerate}

\subsection{Contoh}

Untuk grammar
\[
\begin{aligned}
E &\to T E', &
E' &\to + T E' \mid \varepsilon,\\
T &\to F T', &
T' &\to * F T' \mid \varepsilon,\\
F &\to (E) \mid \code{id},
\end{aligned}
\]
hasil perhitungan adalah:
\begin{table}[H]
\centering
\caption{FIRST dan FOLLOW grammar ekspresi}
\label{tab:ch10-first-follow}
\begin{tabular}{@{}l l l@{}}
\toprule
\textbf{Nonterminal} & \textbf{FIRST} & \textbf{FOLLOW} \\
\midrule
\(E\)  & \(\{ \code{id}, ( \}\) & \(\{ \$, ) \}\) \\
\(E'\) & \(\{ +, \varepsilon \}\) & \(\{ \$, ) \}\) \\
\(T\)  & \(\{ \code{id}, ( \}\) & \(\{ \$, ), + \}\) \\
\(T'\) & \(\{ *, \varepsilon \}\) & \(\{ \$, ), + \}\) \\
\(F\)  & \(\{ \code{id}, ( \}\) & \(\{ \$, ), *, + \}\) \\
\bottomrule
\end{tabular}
\end{table}

\subsection{Alat Bantu}

Perhitungan manual rawan kesalahan. Alat \file{first\_follow.py}
menghitung FIRST dan FOLLOW secara otomatis. Potongan fungsi
\code{compute\_first} ditunjukkan pada \lstref{lst:ch10-first-code}.

\includecode[
  language=Python,
  firstline=61,
  lastline=99,
  firstnumber=61,
  caption={Fungsi \code{compute\_first} pada \file{first\_follow.py}},
  label={lst:ch10-first-code}
]{../../code/tools/first_follow.py}

Alat ini membaca grammar dalam format sederhana dan mencetak
FIRST/FOLLOW. Gunakan alat untuk memeriksa hasil manual sebelum
membangun tabel LL(1).

\subsection{Kondisi LL(1)}

Sebuah grammar bersifat LL(1) jika untuk setiap nonterminal \(A\)
dan setiap pasangan aturan \(A \to \alpha\) dan \(A \to \beta\)
berlaku:
\[
\text{FIRST}(\alpha) \cap \text{FIRST}(\beta) = \varnothing,
\]
dan jika salah satu nullable, maka
\[
\text{FIRST}(\alpha) \cap \text{FOLLOW}(A) = \varnothing.
\]
Jika kondisi ini dilanggar, tabel parsing akan memiliki sel ganda
dan grammar bukan LL(1).
